% Short running form; the full title is too wide for the head's right box.
\runninghead{The Will None Resists}
\section{The Will None Resists, and a Servant Who Would Not Forgive}
\label{sec:opening}

The Mass of the Twenty-first Sunday after Pentecost, \latin{Dominica vigesima
prima post Pentecosten}, opens with a creed spoken under sentence of death. Its Introit is taken
from the prayer of Mardochai in the Greek chapters of the Book of Esther
(Esth 13:9--11), prayed when the king's decree, obtained by Aman, had condemned
the Jews of the Persian empire to destruction. The 1861 Philadelphia hand missal
renders its first sentence ``All things, O Lord, are in thy power: and no one
can resist thy will''; the antiphon goes on to confess God the maker of heaven
and earth and of all that the compass of heaven contains, and ends \latin{Dominus
universorum tu es}. Its verse is the first of the longest psalm, ``Blessed are
the undefiled in the way, who walk in the law of the Lord.'' The Collect asks
God to keep his family by his continual mercy, \latin{continua pietate}, free
under his protection from every adversity and devoted to his name in good works.
The Epistle, from the last chapter of Ephesians, arms the Church: ``our
wrestling is not against flesh and blood; but against principalities and
powers, against the rulers of the world of this darkness,'' and so the
Christian is to take the whole armour of God, ``that you may be able to resist
in the evil day.'' The Gradual answers with the opening of Ps~89: ``Lord, thou
hast been our refuge'' from generation to generation, ``from eternity and to
eternity thou art God.'' The Alleluia sings the first verse of the Exodus psalm,
``When Israel went out of Egypt, the house of Jacob from a barbarous people.''

The Gospel is the parable of the unforgiving servant (Mt 18:23--35). A king
settles accounts with his servants; one owes ten thousand talents, cannot pay,
is ordered to be sold with his wife and children, falls down and begs, ``Have
patience with me, and I will pay thee all,'' and is let go with the whole debt
forgiven. Going out, he meets a fellow servant who owes him a hundred pence,
takes him by the throat, hears the same words, ``Have patience with me, and I
will pay thee all,'' and casts him into prison. The king calls him back: ``Thou
wicked servant, I forgave thee all the debt, because thou besoughtest me:
shouldst not thou then have had compassion also on thy fellow servant, even as
I had compassion on thee?'' He hands him to the torturers until he pays all the
debt. ``So also shall my heavenly Father do to you, if you forgive not every one
his brother from your hearts.'' The Offertory then tells, in a narrative of its
own, how Satan asked to tempt Job, ``and power was given him by the Lord over his
possessions, and over his flesh,'' in the words of the 1861 hand missal. The
Secret offers the sacrifice by which God willed to be appeased and to restore
salvation by his powerful mercy. The Communion takes three half-verses of the
eleventh strophe of Ps~118: the soul that has hoped in God's word asks, ``when
wilt thou judge those who persecute me?'' and begs his help. The Postcommunion
asks that the food of immortality received with the mouth may be followed with
a pure mind.

A sovereign will that no one resists, and an adversary who asks for power and
is given it; a debtor forgiven without measure who will not forgive a little;
a people brought out of Egypt that asks when its persecutors will be judged.
Three questions rise from these texts, and each gives one of the readings that
follow. The first asks who the Christian's enemy is, and how the adversary's
work can stand within a will that none resists. St Gregory the Great answers
from Job that God looses the enemy and bridles him, permitting for the proving
of his servant what the devil asks for his ruin; St John Chrysostom, from the
Epistle, that the spirits who work in the persecutor are the real enemy, so that
the wronged man is to turn his anger on the devil and not on his brother. The
second asks what the king's mercy requires of the forgiven, and whether a pardon
once given can return. Chrysostom answers that the servant revoked the gift
himself by his cruelty; St Augustine, that God calls back what he remitted from
the one who will not forgive, and that forgiveness from the heart leaves room
for correction given in love. The third asks when the judgment that the
Communion asks for will come, and where that leaves the people that sings this
Mass in the meantime. Augustine hears in the plea the martyrs' cry under the
altar, answered with long-suffering until their brethren are complete;
Chrysostom reads the girded loins of the Epistle as the girding of the Passover
eaten on the night of the Exodus, and the Christian life as a march toward
departure. Every element of the formulary has a place in each reading, and each
closes with four senses of its own, literal, allegorical, moral and anagogical.

The Sunday is of the second class in the Proper of Time of the 1962
\work{Missale Romanum}. In 2026 it falls on 18~October, which the calendar gives to St Luke
the Evangelist, a feast of the second class. A Sunday of the second class is
preferred to a feast of the same class that is not a feast of the Lord
(\work{Rubricae generales} 16), and St Luke is commemorated wherever the rubrics
admit an ordinary commemoration: in every low Mass and in a conventual Mass (RG
108, 111~b; \work{Rubricae generales Missalis} 434~b). There his Collect, Secret
and Postcommunion follow the Sunday's, each under its own conclusion (RGMR
437~a), and nothing else of his formulary comes with them, not even the Preface
of the Apostles (RGMR 483). A sung Mass that is not conventual, such as a parish
\latin{Missa cantata}, admits no such commemoration (RG 111~a; RGMR 434~a). The
date carries a second addition. The eighteenth is the penultimate Sunday of
October, the day the general rubrics give to the missions, and on it the
oration \latin{pro Fidei propagatione} is joined to the Collect, the Secret and
the Postcommunion of the Mass, each time under one conclusion with the Sunday's
prayer, in every Mass, sung or low (RGMR 444~c, 445--446, 453), unless the
Ordinary of the place has appointed another day for the missions. Its Collect
begins \latin{Deus, qui omnes homines vis salvos fieri}, the words of 1~Tim 2:4,
and so it is prayed on this Sunday beside an Introit that says no one can resist
God's will. Neither the commemoration nor the added oration belongs to the
formulary studied here. The Mass is said in green (RG 127~b), with the Gloria,
with the Credo the rubrics require on every Sunday (RGMR 475~a), and with the
Preface of the Most Holy Trinity, which the Missal prints as a direction after
the Secret and the general rubrics appoint for the Sundays of the second class
outside the Christmas and Easter seasons (RGMR 494~b).

\subsection{The appointed elements at a glance}

\begin{maptable}
Introit & \latin{In voluntate tua, Domine}; Esth 13:9--11, adapted; verse
Ps 118:1 (Heb.\ 119) &
Mardochai's confession under sentence of death: ``All things, O Lord, are in
thy power: and no one can resist thy will'' (1861 English); God made heaven and
earth and is Lord of all. The verse blesses ``the undefiled in the way, who walk
in the law of the Lord''.\\
Collect & \latin{Familiam tuam, quaesumus, Domine} &
In the 1861 English, ``Preserve thy family, O Lord, we beseech thee, by thy
constant mercy: that under thy protection, it may be freed from all adversity;
and by the practice of good works, continue devoted to thy name.''\\
Epistle & Eph 6:10--17 &
Be strengthened in the Lord; put on the armour of God against the deceits of the
devil, for the wrestling is not against flesh and blood but against
principalities, powers and spirits of wickedness; stand in the evil day girt
with truth, with the breastplate of justice, the gospel of peace, the shield of
faith, the helmet of salvation and the sword of the Spirit, ``which is the word
of God''.\\
Gradual & \latin{Domine, refugium}; Ps 89:1--2 (Heb.\ 90) &
``Lord, thou hast been our refuge from generation to generation. Before the
mountains were made, or the earth and the world was formed; from eternity and to
eternity thou art God.'' The psalm's title, naming Moses, is not sung.\\
Alleluia & \latin{In exitu Israel}; Ps 113:1 (Heb.\ 114:1) &
``When Israel went out of Egypt, the house of Jacob from a barbarous people.''\\
Gospel & Mt 18:23--35 &
The king who settles accounts; the servant forgiven ten thousand talents who
throttles a fellow servant for a hundred pence; the king's sentence, ``So also
shall my heavenly Father do to you, if you forgive not every one his brother
from your hearts.''\\
Offertory & \latin{Vir erat in terra Hus}; printed citation Job 1 &
A narrative of Job's testing: a man simple and upright and fearing God, ``and
Satan asked to tempt him; and power was given him by the Lord over his
possessions, and over his flesh'' (1861 English); he lost his substance and his
sons and was struck with a grievous ulcer.\\
Secret & \latin{Suscipe, Domine, propitius hostias} &
In the 1861 English, ``Mercifully receive, O Lord, these offerings, by which
thou art pleased to be appeased, and restore us to salvation by thy powerful
goodness.''\\
Communion & \latin{In salutari tuo}; Ps 118:81, 84 and 86 (Heb.\ 119), adapted &
The soul has hoped in God's salvation and his word; ``when wilt thou judge those
who persecute me?''; the wicked have persecuted me, ``help me, O Lord my God''
(1861 English). The psalm's question about the number of the servant's days is
not sung.\\
Postcommunion & \latin{Immortalitatis alimoniam} &
In the 1861 English, ``Having received the food of immortality, we beseech
thee, O Lord, that what we have taken with our mouths we may receive with pure
souls.''\\
\end{maptable}
